%%%%%%%%%%%%%%%%%%%%%%%%%%%%%%%%%%%%%%%%%%%%%%%%%
\section{Renormalization}
\label{sec:renormalization}

In this section, we present our general arguments for the renormalization of the three topologies of diagrams that could give UV divergent contribution to the quasi-quark PDFs, as identified in the last section.

Let us first consider the power UV divergence from the general diagram in Fig.~\ref{fig:multiloopdiv}(a), which could be expressed in terms of the sum of diagrams made of 1PI diagrams, as shown in Fig.~\ref{fig:sumSE}.  Since the UV divergence is local in coordinate space, we assign the ``i-th'' blob in Fig.~\ref{fig:sumSE} a specific point $r_i$. We express the linear UV power divergence of the ``1st'' blob in Fig.~\ref{fig:sumSE} as $c \int_0^{r} d r_1$, where perturbative coefficient $c=c^{(1)}+c^{(2)}+\cdots$ where $c^{(1)}=-\frac{\alpha_s C_F}{\pi}\frac{1}{a} $ is the first order expansion in $\alpha_s$, as derived in Eq.~(\ref{eq:1a}).

With the geometric sum of the 1PI diagrams, as shown in Fig.~\ref{fig:sumSE}, we can derive all linear UV power divergent contribution to the general diagram in Fig.~\ref{fig:multiloopdiv}(a) by summing over contributions of 1PI diagrams located between 0 and $r\,(r>0)$ to all orders as,
\begin{eqnarray}
&\ & 1+ c \int_0^{r} d r_1 + c^2 \int_0^{r} d r_1 \int_{r_1}^{r} d r_2 + \cdots
\nonumber \\
&\ & \hskip 0.3in
= {\cal P} e^{\, c \int_0^r dr'} = e^{\, c\, r}\, ,
\end{eqnarray}
with ${\cal P}$ indicating the path ordering.
Therefore, one could introduce an overall factor $e^{-c\, |\xi_z|}$ to remove all linear power UV divergences of the quasi-quark PDFs. This overall factor could be thought as the mass renormalization of a test particle moving along the gauge link. Renormalizing power divergence in this way was first proposed in Ref.~\cite{Polyakov:1980ca}.

Besides power UV divergence, there are also logarithmic UV divergences from Fig.~\ref{fig:multiloopdiv}(a).  It is well known~\cite{Dotsenko:1979wb} that these divergences can be removed by a ``wave function'' renormalization of the test particle, $Z_{wq}^{-1}$.

The general diagrams in Fig.~\ref{fig:multiloopdiv}(b) have only the logarithmic UV divergences, effectively from the loop corrections to the gluon coupling to the gauge link.  it was proven in Ref.~\cite{Dotsenko:1979wb} that these logarithmic UV divergences can be absorbed by the coupling constant renormalization of QCD. Therefore, we do not need to worry about them if we use the renormalized QCD Lagrangian.

Finally, let us examine UV divergence from the general diagrams shown in Fig.~\ref{fig:multiloopdiv}(c). Unlike the general diagrams in Figs.~\ref{fig:multiloopdiv}(a) and (b), the loop momentum of diagrams in Fig.~\ref{fig:multiloopdiv}(c) go through active quark (or gluon for the case of quasi-gluon PDF).  Since the UV divergence from the diagrams in Fig.~\ref{fig:multiloopdiv}(c) effectively comes from high order loop corrections to the quark-gauge-link vertex, which is not a fundamental coupling in QCD Lagrangian, using the renormalized QCD Lagrangian to do the calculation does not help remove this kind of UV divergences.  That is, the renormalziation of the operators defining quasi-PDFs is required to remove this kind of UV divergences, if we want to have renormalizable quasi-PDFs.

The key question is then if the operators defining quasi-PDFs will mix with other operators under renormalization? If it does, there could be a danger that the operators defining quasi-PDFs may not form a closed set under the renormalization.

The quark-gauge-link vertex at the lowest order is a simple gamma matrix $\gamma_z$ for the quasi-quark PDFs (is the same for the $A_z=0$ case). As we demonstrated in the last section, UV divergence comes only from the region where all loop momenta are very large, thus we can set $p=0$ if we are only interested in leading UV divergence, which is logarithmic as demonstrated in Sec.~\ref{sec:powercounting}. In addition, the logarithmic UV divergence, $\ln(a)$ as $a\to 0$, is local in coordinate space, and does not have a direct dependence on $\xi_z$.  That is, we find that the UV divergent term of Fig.~\ref{fig:multiloopdiv}(c) only depends on the vector $n_z^\mu$, and consequently, it is proportional to $\gamma_z$ with a constant coefficient, which is proportional to the quark-gauge-link vertex at the lowest order.  Therefore, a constant counterterm is sufficient to remove this kind of UV divergences. Using bookkeeping forests subtraction method, it is straight forward to remove the high order divergences and to show that the net effect is to introduce a constant multiplicative renormalizaton factor $Z_{vq}^{-1}$ for the quark-gauge-link vertex.

In summary, by using renormalized QCD Lagrangian in Feynman gauge, we find that all remained perturbative UV divergences of the quasi-quark PDFs can be removed by introducing a multiplicative renormalization factor, with the multiplicative renormalization factor calculated order by order in QCD perturbation theory.  We expect that similar arguments should also apply for quasi-gluon PDF.  We thus can define the renormalized coordinate-space quasi-PDFs as
\begin{align}\label{eq:renor}
\begin{split}
&\tilde{F}^{R}_{i/p}(\xi_z,\tilde{\mu}^2,p_z)=e^{-C_i |\xi_z|}Z_{wi}^{-1}Z_{vi}^{-1}\tilde{F}^{b}_{i/p}(\xi_z,\tilde{\mu}^2,p_z),
\end{split}
\end{align}
where $C_i$, $Z_{wi}$ and $Z_{vi}$ are renormalization constant depending on parton flavor ``$i$'' but independent of $\xi_z$, and perturbatively calculable order-by-order in powers of $\alpha_s$.   We conclude that the coordinate-space quasi-PDFs are renormalizable and they do not mix with each other or with any other operators under renormalization group equation.
